\section{Repetitions}
%\label{sec:7:1:repetitions}

We have already seen a few examples of \python\ programs in which a
sequence of \robot\ commands or of procedure calls is repeated
many times in a row. This is the case, for example, in the program {\idt FriezeOfN\_1}, whose execution produces the frieze below:

%\vspace{\baselineskip}
\vspace*{8mm}%
\noindent\includevspacePostScript[1cm]{FriseDeN}{0.95}{0.95}
\vspace*{-1cm}%
{
\parindent \parindentprog
\parskip 0pt
\nopagebreak\medskip

\carrenoir
\vspace{\neglskip}\mbox{\blc {\mcl from}{\idt \ LibRobot\ }{\mcl import}{\idt \ *}}\par
\nopagebreak
\vspace{0.3\baselineskip}
\vspace{\neglskip}\mbox{\blc {\cmt \#\ Draw\ a\ frieze\ of\ N}}\par
\vspace{\neglskip}\mbox{\blc {\idt gl()}}\par
\vspace{\neglskip}\mbox{\blc {\idt fd();\ fd();\ ht();\ cl();\ cr();\ ht()}}\par
\vspace{\neglskip}\mbox{\blc {\idt fd();\ fd();\ ht();\ cl();\ cr();\ ht()}}\par
\vspace{\neglskip}\mbox{\blc {\idt fd();\ fd();\ ht();\ cl();\ cr();\ ht()}}\par
\vspace{\neglskip}\mbox{\blc {\idt fd();\ fd();\ ht();\ cl();\ cr();\ ht()}}\par
\vspace{\neglskip}\mbox{\blc {\idt fd();\ fd();\ ht();\ cl();\ cr();\ ht()}}\par
\vspace{\neglskip}\mbox{\blc {\idt fd();\ fd();\ ht();\ cl();\ cr();\ ht()}}\par
\vspace{\neglskip}\mbox{\blc {\idt fd();\ fd();\ ht();\ cl();\ cr();\ ht()}}\par
\vspace{\neglskip}\mbox{\blc {\idt fd();\ fd();\ ht();\ cl();\ cr();\ ht()}}\par
\vspace{\neglskip}\mbox{\blc {\idt fd();\ fd();\ ht();\ cl();\ cr();\ ht()}}\par
\vspace{\neglskip}\mbox{\blc {\idt fd();\ fd();\ ht();\ cl();\ cr();\ ht()}}\par
\vspace{\neglskip}\mbox{\blc {\idt fd();\ fd();\ ht();\ cl();\ cr();\ ht()}}\par
\vspace{\neglskip}\mbox{\blc {\idt fd();\ fd();\ ht();\ cl();\ cr();\ ht()}}\par
\vspace{\neglskip}\mbox{\blc {\idt fd();\ fd();\ ht();\ cl();\ cr();\ ht()}}\par
\vspace{\neglskip}\mbox{\blc {\idt fd();\ fd();\ ht();\ cl();\ cr();\ ht()}}\par
\vspace{\neglskip}\mbox{\blc {\idt fd();\ fd();\ ht();\ cl();\ cr();\ ht()}}\par
\vspace{\neglskip}\mbox{\blc {\idt fd();\ fd();\ ht();\ cl();\ cr();\ ht()}}\par
\nopagebreak
\vspace{\neglskip}\mbox{\blc {\idt st()}}\par

\nopagebreak\medskip }\nopagebreak


To avoid repeating the same sequences of in\-struc\-tions over and over,
we are going to study a new statement of the \python\ language called the
{\em ``for'' loop}\index{For loop@``for'' loop} (``for'' is an ordinary English word, so the statement reads almost like an English sentence%
\footnote{\ We call it a ``for'' loop, using the very word that is
written in the program.}). Here is the
program {\idt FriezeOfN\_1} drawing the frieze below,
rewritten using a ``for'' loop:

%\vspace{\baselineskip}
\vspace*{8mm}%
\noindent\includevspacePostScript[1cm]{FriseDeNprime}{0.95}{0.95}
\vspace*{-1cm}%
\label{prog:FriseDeN_2}
{
\parindent \parindentprog
\parskip 0pt
\nopagebreak\medskip

\carrenoir
\vspace{\neglskip}\mbox{\blc {\mcl from}{\idt \ LibRobot\ }{\mcl import}{\idt \ *}}\par
\nopagebreak
\vspace{0.3\baselineskip}
\vspace{\neglskip}\mbox{\blc {\cmt \#\ Draw\ a\ frieze\ of\ N\ with\ a\ \char34{}for\char34{}\ loop}}\par
\nopagebreak
\vspace{\neglskip}\mbox{\blc {\idt gl()}}\par
\nopagebreak
\vspace{\neglskip}\mbox{\blc {\mcl for}{\idt \ j\ }{\mcl in}{\idt \ range(16):}}\par
\nopagebreak
\vspace{\neglskip}\mbox{\blc \blk\blk{\idt fd();\ fd();\ ht();\ cl();\ cr();\ ht()}}\par
\nopagebreak
\vspace{\neglskip}\mbox{\blc {\idt st()}}\par

\nopagebreak\medskip }\nopagebreak


\noindent In plain English, the loop `{\mcl for}\ {\idt j}\ {\mcl in}\ {\idt range(16):}\ \ldots'
can be read as `for j among the first 16 numbers do \ \ldots'. The English word {\mcl in}
has its usual meaning. The English word {\idt range} means
`sequence' or `span'.

The sequence of statements to be repeated is written on the lines that follow the colon `{\idt :}',
shifted to the right. It is called the
{\em body of the loop}\index{Loop body}. The \index{Iteration}{\em number of iterations}, that is, the
number of times the body of the loop must be repeated, is given
between the parentheses of {\idt range}. In our example, the body of the loop
`{\idt  fd();\ fd();\ ht();\ cl();\ cr();\ ht()}' must be repeated 16 times
in a row.

\section{The loop counter}
%\label{sec:7:3:decls}

Just after the word {\mcl for} we find the name {\idt j} of a
{\em counter}\index{Loop counter}. The computer uses this counter to
count the number of times it has to repeat the sequence of statements in the
``for'' loop, exactly as we might do by counting on our fingers!
To do so, the computer uses a cell of its memory, which successively contains
the integers 0, 1, 2, \ldots, 15. Indeed, {\idt range(16)} denotes the sequence of
16 integers that starts at 0 and stops just before 16 (like computer scientists,
\python\ counts from 0). We say that {\idt j} is an integer
{\em variable}\index{Variable}\index{Integer variable} since its value varies
from one iteration to the next.

We have to give a name to this counter. It is an identifier that we can choose
as we please (here, for example, {\idt  j}), following the rule for
identifiers: it begins with a letter (or an underscore `\_'), contains only letters, digits and underscores, and must not be a reserved word of
\python. If we write `{\mcl for}\ {\idt 1st}\ {\mcl in}\ \ldots', \python\ answers
`{\tt SyntaxError: invalid decimal literal}' (since
{\idt 1st} begins with a digit). If we write `{\mcl for}\ {\mcl in}\ {\mcl in}\ \ldots',
it answers `{\tt SyntaxError: invalid syntax}'.

In \python, there is no need to announce in advance that we are going to use
a counter%
\footnote{\ Other languages such as \pascal\ require the counter to be {\em declared} at the beginning
of the program or of the procedure, stating that it denotes an integer.}:
the counter is created by the ``for'' loop itself. \python\
knows that it is an integer since it takes the values 0, 1, 2, \ldots\ After the
loop, the counter keeps the last value it took (15 in our example), but
it is better not to use it: if the loop did not perform any iteration, the
counter does not exist and \python\ reports the error
`{\tt NameError: name 'j' is not defined}'.

For the counter to take the values 1, 2, 3, \ldots, 16, we would write
`{\mcl for}\ {\idt j}\ {\mcl in}\ {\idt range(1,\ 17):}'. In general,
{\idt range(}$a${\idt ,\ }$b${\idt )} denotes the sequence of integers from $a$ to $b-1$:
it starts at $a$ and stops just {\em before} $b$. Thus {\idt range(1,\ 5)} denotes
1, 2, 3, 4 and {\idt range(}$n${\idt )}, which is an abbreviation of
{\idt range(0,\ }$n${\idt )}, denotes the $n$ numbers 0, 1, \ldots, $n-1$.
For the counter to go from 1 to $n$, we therefore write
`{\mcl for}\ {\idt i}\ {\mcl in}\ {\idt range(1,\ }$n${\idt \ +\ 1):}'.

The loop counter of the program {\idt FriezeOfN\_2} is only used to count the iterations.
It is then simpler to write {\idt range(16)}.
We will see in the next chapter that its value can also be used in the body of
the loop.

\section{The ``for'' loop}
%\label{sec:7:4:for}

A {\em ``for'' loop}\index{For loop@``for'' loop} has the following form:\\
\begin{center}
\begin{tabular}{@{}|p{12.5cm}|@{}} \hline
{\em ``For'' loop:}\\ \hline \hline
\begin{minipage}[t]{12.5cm}
{\medskip
\mbox{\blc\ \ \ {\mcl for}\hspace{0.3em}\blanc\hspace{0.3em}{\small\rm Identifier}\hspace{0.3em}\blanc\hspace{0.3em}{\mcl in}\hspace{0.3em}\blanc\hspace{0.3em}{\idt range(}{\small\rm Number of iterations}{\idt ):}}\par
\mbox{\blc\ \ \ \ \ \ \ \small\rm  \ldots\ Statements of the body of the loop (indented) \ldots}\par
\medskip}
\end{minipage}\\ \hline
\end{tabular}
\end{center}
\noindent The identifier is the name of the iteration counter.
The sequence of statements that is repeated is called the {\em body of
the loop}\index{Loop body}. Like the body of a procedure, it is written on the lines
that follow the colon `{\idt :}', all shifted to the right by the same number of
blanks (usually four), one statement per line or several statements on the same
line separated by `{\idt ;}'. It is this shift, called {\em indentation}\index{Indentation},
that tells \python\ where the body of the loop begins and ends: the loop
ends at the first line that is no longer shifted.

Do not forget the blanks \hspace{0.3em}\blanc\hspace{0.3em} around the word {\mcl in}.
If we write {\idt jin} instead of {\idt j}\ {\mcl in}, or {\idt inrange} instead of
{\mcl in}\ {\idt range}, \python\ no longer recognizes the word {\mcl in} and reports
`{\tt SyntaxError: invalid syntax}'. If we forget the colon `{\idt :}' at the end
of the {\mcl for} line, \python\ reports `{\tt SyntaxError: expected ':'}'.

A ``for'' loop is itself a statement. The statements that follow it
are written without any shift, at the same level as the {\mcl for}, like the {\idt st()} of the
program {\idt FriezeOfN\_2}. If we also shifted the line {\idt st()}, it would be part
of the body of the loop and would be repeated 16 times! If we forget to shift the body
of the loop, \python\ reports
`{\tt IndentationError: expected an indented block after 'for' statement on line 5}'.
If the lines of the body are not all shifted by the same number of blanks,
\python\ reports `{\tt IndentationError: unexpected indent}' or
`{\tt IndentationError: unindent does not match any outer indentation level}'.

If the `Number of iterations' is smaller than 1 (0, for example), the statements of the
body of the loop are not executed at all. \underline{Do not put} \underline{a semicolon `{\idt $;$}'} \underline{after the colon}: \python\ then reports `{\tt SyntaxError: invalid syntax}' just after the colon.

When the body of the loop is short, the whole loop can be written on a single
line, as in `{\mcl for}\ {\idt i\ }{\mcl in}\ {\idt range(10):\ fd()}'.
Extra blanks inside a line only serve to make the
program easier to read, but those at the beginning of a line matter.

\begin{exemple}[Frieze of spikes]{\rm
 The following frieze:\\[20mm]
\includevspacePostScript[1cm]{FriseDePointes}{0.95}{0.95}%
\\[1ex]
is produced by the program below:
{
\parindent \parindentprog
\parskip 0pt
\nopagebreak\medskip

\carrenoir
\vspace{\neglskip}\mbox{\blc {\mcl from}{\idt \ LibRobot\ }{\mcl import}{\idt \ *}}\par
\nopagebreak
\vspace{0.3\baselineskip}
\vspace{\neglskip}\mbox{\blc {\cmt \#\ Draw\ a\ frieze\ of\ spikes}}\par
\nopagebreak
\vspace{\neglskip}\mbox{\blc {\idt gl();\ qtr()}}\par
\nopagebreak
\vspace{\neglskip}\mbox{\blc {\mcl for}{\idt \ j\ }{\mcl in}{\idt \ range(32):}}\par
\nopagebreak
\vspace{\neglskip}\mbox{\blc \blk\blk{\idt pu();\ fd();\ pd();\ ht();\ cr();\ ht()}}\par
\nopagebreak
\vspace{\neglskip}\mbox{\blc \blk\blk{\idt fd();\ fd();\ ht();\ cr()}}\par
\nopagebreak
\vspace{\neglskip}\mbox{\blc {\idt st()}}\par

\nopagebreak\medskip }\nopagebreak

\noindent $\Box$}
\end{exemple}

\section{Using a ``for'' loop in a procedure}
%\label{sec:7:5:for-proc}

The body of a procedure can contain ``for'' loops, like any
other statement:
\index{Loop counter}
\begin{center}
\begin{tabular}{@{}|p{12.5cm}|@{}} \hline
{\em Procedure declaration:}\\ \hline \hline
\begin{minipage}[t]{12.5cm}
\medskip
\mbox{\blc\ \ \ {\mcl def}\hspace{0.3em}\blanc\hspace{0.3em}\idt procedureName():}\par
\mbox{\blc\ \ \ \ \ \ \ \small\rm \ldots\ Statements of the body of the procedure%
\footnote{\ Including ``for'' loops using counters {\idt i}, {\idt j}, \ldots}
(indented) \ldots}\par
\medskip
\end{minipage}\\ \hline
\end{tabular}
\end{center}

Since the counter of a loop is only used in the body of that loop and since it
starts again from its first value at each new loop, the same
identifier can serve as a counter in several different loops,
whether they are in different procedures or one after the other
in the same procedure. This is illustrated in the program
{\idt Arrow} below, where different loop
counters, used in different procedures ({\idt
fletching} and {\idt shaft}), have the same name {\idt i}.

\noindent\begin{minipage}{\textwidth}
\begin{exemple}[Arrow]{\rm \ \\[7cm]
\includePostScript{FlecheVerticale}{0.95}{0.95}%
\vspace*{-7cm}%
{
\parindent \parindentprog
\parskip 0pt
\nopagebreak\medskip

\carrenoir
\vspace{\neglskip}\mbox{\blc {\mcl from}{\idt \ LibRobot\ }{\mcl import}{\idt \ *}}\par
\nopagebreak
\vspace{0.3\baselineskip}
\vspace{\neglskip}\mbox{\blc {\cmt \#\ Draw\ a\ vertical\ arrow}}\par
\vspace{\neglskip}\mbox{\blc {\mcl def}{\idt \ feather():}}\par
\vspace{\neglskip}\mbox{\blc \blk\blk{\idt fd();\ tetl();\ fd();\ pu();\ tetl();\ fd();\ fd()}}\par
\vspace{\neglskip}\mbox{\blc \blk\blk{\idt tetl();\ pd();\ fd();\ etr()}}\par
\vspace{0.3\baselineskip}
\vspace{\neglskip}\mbox{\blc {\mcl def}{\idt \ fletching():}}\par
\vspace{\neglskip}\mbox{\blc \blk\blk{\mcl for}{\idt \ i\ }{\mcl in}{\idt \ range(5):}}\par
\vspace{\neglskip}\mbox{\blc \blk\blk\blk\blk{\idt feather()}}\par
\vspace{0.3\baselineskip}
\vspace{\neglskip}\mbox{\blc {\mcl def}{\idt \ shaft():}}\par
\vspace{\neglskip}\mbox{\blc \blk\blk{\mcl for}{\idt \ i\ }{\mcl in}{\idt \ range(10):}}\par
\vspace{\neglskip}\mbox{\blc \blk\blk\blk\blk{\idt fd()}}\par
\vspace{0.3\baselineskip}
\vspace{\neglskip}\mbox{\blc {\mcl def}{\idt \ head():}}\par
\vspace{\neglskip}\mbox{\blc \blk\blk{\idt qtl();\ fd();\ tetr();\ fd();\ qtr();\ fd()}}\par
\vspace{\neglskip}\mbox{\blc \blk\blk{\idt tetr();\ fd()}}\par
\nopagebreak
\vspace{0.3\baselineskip}
\vspace{\neglskip}\mbox{\blc {\idt gb();\ fletching();\ shaft();\ head();\ st()}}\par

\nopagebreak\medskip }\nopagebreak

\noindent $\Box$ }\end{exemple}
\end{minipage}


\section{Nested ``for'' loops (in\-de\-pen\-dent)}
%\label{sec:7:6:for-for-ind}

The body of a ``for'' loop:

\mbox{\blc\ \ \ {\mcl for}\hspace{0.3em}\blanc\hspace{0.3em}{\small\rm Identifier1}\hspace{0.3em}\blanc\hspace{0.3em}{\mcl in}\hspace{0.3em}\blanc\hspace{0.3em}{\idt range(}{\small\rm Number of iterations1}{\idt ):}}\par
\mbox{\blc\ \ \ \ \ \ \ \small\rm  \ldots}\par

\noindent can itself contain ``for'' loops:

\mbox{\blc\ \ \ {\mcl for}\hspace{0.3em}\blanc\hspace{0.3em}{\small\rm Identifier1}\hspace{0.3em}\blanc\hspace{0.3em}{\mcl in}\hspace{0.3em}\blanc\hspace{0.3em}{\idt range(}{\small\rm Number of iterations1}{\idt ):}}\par
\mbox{\blc\ \ \ \ \ \ \ \small\rm  \ldots}\par
\mbox{\blc\ \ \ \ \ \ \ {\mcl for}\hspace{0.3em}\blanc\hspace{0.3em}{\small\rm Identifier2}\hspace{0.3em}\blanc\hspace{0.3em}{\mcl in}\hspace{0.3em}\blanc\hspace{0.3em}{\idt range(}{\small\rm Number of iterations2}{\idt ):}}\par
\mbox{\blc\ \ \ \ \ \ \ \ \ \ \ \small\rm  \ldots}\par
\mbox{\blc\ \ \ \ \ \ \ \small\rm  \ldots}\par

\noindent We then say that the ``for'' loops are {\em nested}.
The loop using the
counter {\small\rm Identifier2} to perform
{\small\rm Number of iterations2} iterations is called the
 {\em inner loop}. The loop using the counter {\small\rm Identifier1} to perform
{\small\rm Number of iterations1} iterations is called the {\em outer loop}.
The body of the inner loop is shifted twice: once because it is part
of the body of the outer loop, a second time because it is the body of the inner
loop. The inner loop ends at the first line that returns to the shift
of its {\mcl for}; the outer loop at the first line that returns to the shift of the
first {\mcl for}. So it is the shift that indicates, without ambiguity, which {\mcl for}
each statement belongs to. Consider, for example,
the drawing of a square:
\par\noindent\begin{minipage}{\linewidth}
\vbox to 0pt{\includevspacePostScript[20mm]{Carre1}{0.95}{0.95}}%
{
\parindent \parindentprog
\parskip 0pt
\nopagebreak\medskip

\carrenoir
\vspace{\neglskip}\mbox{\blc {\mcl from}{\idt \ LibRobot\ }{\mcl import}{\idt \ *}}\par
\nopagebreak
\vspace{0.3\baselineskip}
\vspace{\neglskip}\mbox{\blc {\cmt \#\ Draw\ a\ square\ with\ procedures}}\par
\nopagebreak
\vspace{0.3\baselineskip}
\vspace{\neglskip}\mbox{\blc {\cmt \#\ Draw\ one\ side\ of\ the\ square\ and\ turn\ right}}\par
\nopagebreak
\vspace{\neglskip}\mbox{\blc {\mcl def}{\idt \ side():}}\par
\nopagebreak
\vspace{\neglskip}\mbox{\blc \blk\blk{\mcl for}{\idt \ j\ }{\mcl in}{\idt \ range(4):}}\par
\nopagebreak
\vspace{\neglskip}\mbox{\blc \blk\blk\blk\blk{\idt fd()}}\par
\nopagebreak
\vspace{\neglskip}\mbox{\blc \blk\blk{\idt qtr()}}\par
\nopagebreak
\vspace{0.3\baselineskip}
\vspace{\neglskip}\mbox{\blc {\cmt \#\ Draw\ a\ square}}\par
\nopagebreak
\vspace{\neglskip}\mbox{\blc {\mcl def}{\idt \ square():}}\par
\nopagebreak
\vspace{\neglskip}\mbox{\blc \blk\blk{\mcl for}{\idt \ i\ }{\mcl in}{\idt \ range(4):}}\par
\nopagebreak
\vspace{\neglskip}\mbox{\blc \blk\blk\blk\blk{\idt side()}}\par
\nopagebreak
\vspace{0.3\baselineskip}
\vspace{\neglskip}\mbox{\blc {\idt square();\ st()}}\par

\nopagebreak\medskip }\nopagebreak

\end{minipage}

By replacing the procedure calls by their bodies, we obtain
two nested loops:

%\includePostScript{Carre1}{0.95}{0.95}
{
\parindent \parindentprog
\parskip 0pt
\nopagebreak\medskip

\carrenoir
\vspace{\neglskip}\mbox{\blc {\mcl from}{\idt \ LibRobot\ }{\mcl import}{\idt \ *}}\par
\nopagebreak
\vspace{0.3\baselineskip}
\vspace{\neglskip}\mbox{\blc {\cmt \#\ Draw\ a\ square\ with\ nested\ \char34{}for\char34{}\ loops}}\par
\nopagebreak
\vspace{\neglskip}\mbox{\blc {\cmt \#\ Draw\ a\ square}}\par
\nopagebreak
\vspace{\neglskip}\mbox{\blc {\mcl for}{\idt \ i\ }{\mcl in}{\idt \ range(4):}}\par
\nopagebreak
\vspace{\neglskip}\mbox{\blc \blk\blk{\cmt \#\ Draw\ one\ side\ of\ the\ square\ and\ turn\ right}}\par
\nopagebreak
\vspace{\neglskip}\mbox{\blc \blk\blk{\mcl for}{\idt \ j\ }{\mcl in}{\idt \ range(4):}}\par
\nopagebreak
\vspace{\neglskip}\mbox{\blc \blk\blk\blk\blk{\idt fd()}}\par
\nopagebreak
\vspace{\neglskip}\mbox{\blc \blk\blk{\idt qtr()}}\par
\nopagebreak
\vspace{\neglskip}\mbox{\blc {\idt st()}}\par

\nopagebreak\medskip }\nopagebreak


The inner loop is used to build one side while the
outer loop draws the four sides.

\newpage
\section{Nested ``for'' loops (dependent)}
%\label{sec:7:7:for-for-dep}

In nested ``for'' loops, the number of iterations
of the in\-ner loop may be dif\-fer\-ent each time the body of
the outer loop is executed. Let us begin by considering the simple
example of line segments of lengths 1, 2, 3, \ldots\ arranged in a spiral
(the drawing is shown below at scale $\frac{1}{2}$ with the grid):\\[5cm]
\includePostScript{SpiraleCarree-1}{0.95}{0.95}%
\vspace*{-5cm}%
{
\parindent \parindentprog
\parskip 0pt
\nopagebreak\medskip

\carrenoir
\vspace{\neglskip}\mbox{\blc {\mcl from}{\idt \ LibRobot\ }{\mcl import}{\idt \ *}}\par
\nopagebreak
\vspace{0.3\baselineskip}
\vspace{\neglskip}\mbox{\blc {\cmt \#\ Draw\ a\ square\ spiral}}\par
\nopagebreak
\vspace{\neglskip}\mbox{\blc {\mcl for}{\idt \ i\ }{\mcl in}{\idt \ range(1,\ 23):}}\par
\nopagebreak
\vspace{\neglskip}\mbox{\blc \blk\blk{\mcl for}{\idt \ j\ }{\mcl in}{\idt \ range(i):}}\par
\nopagebreak
\vspace{\neglskip}\mbox{\blc \blk\blk\blk\blk{\idt fd()}}\par
\nopagebreak
\vspace{\neglskip}\mbox{\blc \blk\blk{\idt qtr()}}\par
\nopagebreak
\vspace{\neglskip}\mbox{\blc {\idt st()}}\par

\nopagebreak\medskip }\nopagebreak

\vspace{2.6cm}% le listing Python est plus court que la figure placee a cote

The number of iterations in the inner ``for'' loop is equal to the value
of the counter {\idt i} of the outer loop, that is, successively 1, 2,
3, \ldots, 22 (since {\idt range(1,\ 23)} stops just before 23). The program {\idt SquareSpiral\_1} is therefore equivalent to the
program {\idt SquareSpiral\_2} below, in which the ``for'' loops have been
eliminated to show exactly which sequence of commands is executed
by the \robot.

%\special{SpiraleCarree-1.ps 0.24 0.24 scale 0 -342 translate}
{
\parindent \parindentprog
\parskip 0pt
\nopagebreak\medskip

\carrenoir
\vspace{\neglskip}\mbox{\blc {\mcl from}{\idt \ LibRobot\ }{\mcl import}{\idt \ *}}\par
\nopagebreak
\vspace{0.3\baselineskip}
\vspace{\neglskip}\mbox{\blc {\cmt \#\ Draw\ a\ square\ spiral\ without\ \char34{}for\char34{}\ loops}}\par
\vspace{\neglskip}\mbox{\blc {\idt fd();\ qtr()}}\par
\vspace{\neglskip}\mbox{\blc {\idt fd();\ fd();\ qtr()}}\par
\vspace{\neglskip}\mbox{\blc {\idt fd();\ fd();\ fd();\ qtr()}}\par
\vspace{\neglskip}\mbox{\blc {\idt fd();\ fd();\ fd();\ fd();\ qtr()}}\par
\vspace{\neglskip}\mbox{\blc {\idt fd();\ fd();\ fd();\ fd();\ fd();\ qtr()}}\par
\vspace{\neglskip}\mbox{\blc {\idt fd();\ fd();\ fd();\ fd();\ fd();\ fd();\ qtr()}}\par
\vspace{\neglskip}\mbox{\blc {\idt fd();\ fd();\ fd();\ fd();\ fd();\ fd();\ fd()}}\par
\vspace{\neglskip}\mbox{\blc {\idt qtr()}}\par
\vspace{\neglskip}\mbox{\blc {\idt fd();\ fd();\ fd();\ fd();\ fd();\ fd();\ fd();\ fd()}}\par
\vspace{\neglskip}\mbox{\blc {\idt qtr()}}\par
\vspace{\neglskip}\mbox{\blc {\idt fd();\ fd();\ fd();\ fd();\ fd();\ fd();\ fd();\ fd()}}\par
\vspace{\neglskip}\mbox{\blc {\idt fd();\ qtr()}}\par
\vspace{\neglskip}\mbox{\blc {\idt fd();\ fd();\ fd();\ fd();\ fd();\ fd();\ fd();\ fd()}}\par
\vspace{\neglskip}\mbox{\blc {\idt fd();\ fd();\ qtr()}}\par
\vspace{\neglskip}\mbox{\blc {\idt fd();\ fd();\ fd();\ fd();\ fd();\ fd();\ fd();\ fd()}}\par
\vspace{\neglskip}\mbox{\blc {\idt fd();\ fd();\ fd();\ qtr()}}\par
\vspace{\neglskip}\mbox{\blc {\idt fd();\ fd();\ fd();\ fd();\ fd();\ fd();\ fd();\ fd()}}\par
\vspace{\neglskip}\mbox{\blc {\idt fd();\ fd();\ fd();\ fd();\ qtr()}}\par
\vspace{\neglskip}\mbox{\blc {\idt fd();\ fd();\ fd();\ fd();\ fd();\ fd();\ fd();\ fd()}}\par
\vspace{\neglskip}\mbox{\blc {\idt fd();\ fd();\ fd();\ fd();\ fd();\ qtr()}}\par
\vspace{\neglskip}\mbox{\blc {\idt fd();\ fd();\ fd();\ fd();\ fd();\ fd();\ fd();\ fd()}}\par
\vspace{\neglskip}\mbox{\blc {\idt fd();\ fd();\ fd();\ fd();\ fd();\ fd();\ qtr()}}\par
\vspace{\neglskip}\mbox{\blc {\idt fd();\ fd();\ fd();\ fd();\ fd();\ fd();\ fd();\ fd()}}\par
\vspace{\neglskip}\mbox{\blc {\idt fd();\ fd();\ fd();\ fd();\ fd();\ fd();\ fd()}}\par
\vspace{\neglskip}\mbox{\blc {\idt qtr()}}\par
\vspace{\neglskip}\mbox{\blc {\idt fd();\ fd();\ fd();\ fd();\ fd();\ fd();\ fd();\ fd()}}\par
\vspace{\neglskip}\mbox{\blc {\idt fd();\ fd();\ fd();\ fd();\ fd();\ fd();\ fd();\ fd()}}\par
\vspace{\neglskip}\mbox{\blc {\idt qtr()}}\par
\vspace{\neglskip}\mbox{\blc {\idt fd();\ fd();\ fd();\ fd();\ fd();\ fd();\ fd();\ fd()}}\par
\vspace{\neglskip}\mbox{\blc {\idt fd();\ fd();\ fd();\ fd();\ fd();\ fd();\ fd();\ fd()}}\par
\vspace{\neglskip}\mbox{\blc {\idt fd();\ qtr()}}\par
\vspace{\neglskip}\mbox{\blc {\idt fd();\ fd();\ fd();\ fd();\ fd();\ fd();\ fd();\ fd()}}\par
\vspace{\neglskip}\mbox{\blc {\idt fd();\ fd();\ fd();\ fd();\ fd();\ fd();\ fd();\ fd()}}\par
\vspace{\neglskip}\mbox{\blc {\idt fd();\ fd();\ qtr()}}\par
\vspace{\neglskip}\mbox{\blc {\idt fd();\ fd();\ fd();\ fd();\ fd();\ fd();\ fd();\ fd()}}\par
\vspace{\neglskip}\mbox{\blc {\idt fd();\ fd();\ fd();\ fd();\ fd();\ fd();\ fd();\ fd()}}\par
\vspace{\neglskip}\mbox{\blc {\idt fd();\ fd();\ fd();\ qtr()}}\par
\vspace{\neglskip}\mbox{\blc {\idt fd();\ fd();\ fd();\ fd();\ fd();\ fd();\ fd();\ fd()}}\par
\vspace{\neglskip}\mbox{\blc {\idt fd();\ fd();\ fd();\ fd();\ fd();\ fd();\ fd();\ fd()}}\par
\vspace{\neglskip}\mbox{\blc {\idt fd();\ fd();\ fd();\ fd();\ qtr()}}\par
\vspace{\neglskip}\mbox{\blc {\idt fd();\ fd();\ fd();\ fd();\ fd();\ fd();\ fd();\ fd()}}\par
\vspace{\neglskip}\mbox{\blc {\idt fd();\ fd();\ fd();\ fd();\ fd();\ fd();\ fd();\ fd()}}\par
\vspace{\neglskip}\mbox{\blc {\idt fd();\ fd();\ fd();\ fd();\ fd();\ qtr()}}\par
\vspace{\neglskip}\mbox{\blc {\idt fd();\ fd();\ fd();\ fd();\ fd();\ fd();\ fd();\ fd()}}\par
\vspace{\neglskip}\mbox{\blc {\idt fd();\ fd();\ fd();\ fd();\ fd();\ fd();\ fd();\ fd()}}\par
\vspace{\neglskip}\mbox{\blc {\idt fd();\ fd();\ fd();\ fd();\ fd();\ fd();\ qtr()}}\par
\nopagebreak
\vspace{\neglskip}\mbox{\blc {\idt st()}}\par

\nopagebreak\medskip }\nopagebreak


\section{Invariants of ``for'' loops}
%\label{sec:7:8:inv}

When we draw a frieze like this one:\\[9mm]
\includevspacePostScript[9mm]{FriseDeTildes}{0.95}{0.95}%
\vspace*{-9mm}%

\noindent with a ``for'' loop, the \robot\ must always be facing
the same direction
and be in the same position with respect to the next pattern to be
drawn by executing
the body of the loop. In other words, at the end of the body of the
loop we must
put the \robot\ back in its starting position:
{
\parindent \parindentprog
\parskip 0pt
\nopagebreak\medskip

\carrenoir
\vspace{\neglskip}\mbox{\blc {\mcl from}{\idt \ LibRobot\ }{\mcl import}{\idt \ *}}\par
\nopagebreak
\vspace{0.3\baselineskip}
\vspace{\neglskip}\mbox{\blc {\cmt \#\ Draw\ a\ frieze\ of\ tildes}}\par
\vspace{0.3\baselineskip}
\vspace{\neglskip}\mbox{\blc {\cmt \#\ Draw\ a\ tilde{-}shaped\ pattern\ \char126{}}}\par
\vspace{\neglskip}\mbox{\blc {\mcl def}{\idt \ tilde():}}\par
\vspace{\neglskip}\mbox{\blc \blk\blk{\idt etr();\ fd();\ etr();\ fd()}}\par
\vspace{\neglskip}\mbox{\blc \blk\blk{\idt qtr();\ fd();\ fd();\ qtl()}}\par
\vspace{\neglskip}\mbox{\blc \blk\blk{\idt fd();\ etl();\ fd()}}\par
\vspace{0.3\baselineskip}
\vspace{\neglskip}\mbox{\blc {\idt gl()}}\par
\vspace{\neglskip}\mbox{\blc {\mcl for}{\idt \ i\ }{\mcl in}{\idt \ range(10):}}\par
\vspace{\neglskip}\mbox{\blc \blk\blk{\idt tilde()}}\par
\vspace{\neglskip}\mbox{\blc \blk\blk{\cmt \#\ Back\ up\ to\ the\ next\ tilde\textquotesingle{}s\ start}}\par
\vspace{\neglskip}\mbox{\blc \blk\blk{\idt pu();\ tetl();\ fd();\ pd()}}\par
\vspace{\neglskip}\mbox{\blc \blk\blk{\cmt \#\ Turn\ the\ robot\ to\ face\ north}}\par
\vspace{\neglskip}\mbox{\blc \blk\blk{\idt qtr()}}\par
\nopagebreak
\vspace{\neglskip}\mbox{\blc {\idt st()}}\par

\nopagebreak\medskip }\nopagebreak

We say that the position of the \robot\ is an {\em invariant}\index{Invariant} of the loop, in the sense
that it must be the same at each iteration.

If, in our example, we draw one pattern and then the next one without putting
the \robot\ back in its starting position in between, we get the following result, which,
it must be admitted, does not look much like the original:\\[5cm]
\includePostScript{FriseDeTildesErr1}{0.95}{0.95}%
\vspace*{-5cm}
{
\parindent \parindentprog
\parskip 0pt
\nopagebreak\medskip

\carrenoir
\vspace{\neglskip}\mbox{\blc {\mcl from}{\idt \ LibRobot\ }{\mcl import}{\idt \ *}}\par
\nopagebreak
\vspace{0.3\baselineskip}
\vspace{\neglskip}\mbox{\blc {\cmt \#\ Draw\ a\ frieze\ of\ tildes\ (first\ mistake)}}\par
\nopagebreak
\vspace{0.3\baselineskip}
\vspace{\neglskip}\mbox{\blc {\cmt \#\ Draw\ a\ tilde{-}shaped\ pattern\ \char126{}}}\par
\nopagebreak
\vspace{\neglskip}\mbox{\blc {\mcl def}{\idt \ tilde():}}\par
\nopagebreak
\vspace{\neglskip}\mbox{\blc \blk\blk{\idt etr();\ fd();\ etr();\ fd()}}\par
\nopagebreak
\vspace{\neglskip}\mbox{\blc \blk\blk{\idt qtr();\ fd();\ fd();\ qtl()}}\par
\nopagebreak
\vspace{\neglskip}\mbox{\blc \blk\blk{\idt fd();\ etl();\ fd()}}\par
\nopagebreak
\vspace{0.3\baselineskip}
\vspace{\neglskip}\mbox{\blc {\mcl for}{\idt \ i\ }{\mcl in}{\idt \ range(10):}}\par
\nopagebreak
\vspace{\neglskip}\mbox{\blc \blk\blk{\idt tilde()}}\par
\nopagebreak
\vspace{\neglskip}\mbox{\blc {\idt st()}}\par

\nopagebreak\medskip }\nopagebreak


If we remember to move the \robot\ back but forget to turn it
back in the right direction, the result is even stranger:

% figure et programme insecables (sinon, a un saut de page, l'espace
% negatif fait remonter le programme au-dessus de l'en-tete)
\noindent\begin{minipage}{\linewidth}
\vspace*{25mm}
\includePostScript{FriseDeTildesErr2}{0.95}{0.95}%
\vspace*{-25mm}
{
\parindent \parindentprog
\parskip 0pt
\nopagebreak\medskip

\carrenoir
\vspace{\neglskip}\mbox{\blc {\mcl from}{\idt \ LibRobot\ }{\mcl import}{\idt \ *}}\par
\nopagebreak
\vspace{0.3\baselineskip}
\vspace{\neglskip}\mbox{\blc {\cmt \#\ Draw\ a\ frieze\ of\ tildes\ (second\ mistake)}}\par
\nopagebreak
\vspace{0.3\baselineskip}
\vspace{\neglskip}\mbox{\blc {\cmt \#\ Draw\ a\ tilde{-}shaped\ pattern\ \char126{}}}\par
\nopagebreak
\vspace{\neglskip}\mbox{\blc {\mcl def}{\idt \ tilde():}}\par
\nopagebreak
\vspace{\neglskip}\mbox{\blc \blk\blk{\idt etr();\ fd();\ etr();\ fd()}}\par
\nopagebreak
\vspace{\neglskip}\mbox{\blc \blk\blk{\idt qtr();\ fd();\ fd();\ qtl()}}\par
\nopagebreak
\vspace{\neglskip}\mbox{\blc \blk\blk{\idt fd();\ etl();\ fd()}}\par
\nopagebreak
\vspace{0.3\baselineskip}
\vspace{\neglskip}\mbox{\blc {\mcl for}{\idt \ i\ }{\mcl in}{\idt \ range(13):}}\par
\nopagebreak
\vspace{\neglskip}\mbox{\blc \blk\blk{\idt tilde()}}\par
\nopagebreak
\vspace{\neglskip}\mbox{\blc \blk\blk{\cmt \#\ Back\ up\ to\ the\ next\ tilde\textquotesingle{}s\ start}}\par
\nopagebreak
\vspace{\neglskip}\mbox{\blc \blk\blk{\idt pu();\ tetl();\ fd();\ pd()}}\par
\nopagebreak
\vspace{\neglskip}\mbox{\blc {\idt st()}}\par

\nopagebreak\medskip }\nopagebreak

\end{minipage}

A good way to find the errors in a ``for'' loop is therefore
to look for the invariant, that is, the position and orientation in
which the \robot\ must be before drawing the pattern by
executing the body of the loop. We must then check that the \robot\ is
correctly placed when we enter the loop for the first time
and that it is correctly placed again after executing the body of the loop.

The same notion of invariant applies to procedures. Before each
call of a procedure, the robot must be in one or several
possible states, which must be perfectly defined by whether the
\robot's pen is up or down
and by the position and orientation
of the \robot\ with respect to the figure drawn by the procedure.

\needspace{6\baselineskip}\noindent \rule{137mm}{0.1mm}\par\nopagebreak %---------------------------------------------
\begin{exercices}{\ \\

\vspace{-\baselineskip}
\tiret Write \python\ programs using ``for'' loops to draw
some of the friezes of figure
\ref{fig:Frise-1}  page \pageref{fig:Frise-1}.

\tiret Write \python\ programs using ``for'' loops in pro\-ce\-dures
to draw some of the friezes of figure
\ref{fig:Frise-2}  page \pageref{fig:Frise-2}.

\tiret Write \python\ programs to draw some of the rosettes,
friezes and tiling of figure
\ref{fig:Ro-Fr-Pa}  page \pageref{fig:Ro-Fr-Pa}, using independent
nested ``for'' loops.

\tiret Write \python\ programs to draw some of the spirals,
grids and other organ pipes of figure
\ref{fig:SpirQuadTuy}  below, using dependent nested
``for'' loops.

\tiret Write a \python\ program to draw the tiling of
figure \ref{fig:PavCLC} below, stating precisely the
invariants of the procedures and of the ``for'' loops.
\noindent $\Box$}\end{exercices}
\begin{figure}[htbp]
 \includevspacePostScript[4.5cm]{SpirQuadTuy}{0.95}{0.95}
 \caption{Grids, squares, spiral and organ pipes}
 \label{fig:SpirQuadTuy}
\end{figure}
\begin{figure}[htbp]
 \includevspacePostScript[7.1cm]{PavageCarresTrianglesCercles}{0.95}{0.95}
 \caption{Tiling of squares, triangles and circles}\label{fig:PavCLC}
\end{figure}
\begin{figure}[htbp]
 \includevspacePostScript[17.8cm]{Frise-1}{0.95}{0.95}
 \caption{Friezes (1)}\label{fig:Frise-1}
\end{figure}
\begin{figure}[htbp]
 \includevspacePostScript[17.8cm]{Frise-2}{0.95}{0.95}
 \caption{Friezes (2)}\label{fig:Frise-2}
\end{figure}
\begin{figure}[htbp]
 \includevspacePostScript[17.8cm]{Ro-Fr-Pa}{0.95}{0.95}
 \caption{Rosettes, friezes and tiling}\label{fig:Ro-Fr-Pa}
\end{figure}

\clearpage

%%%%%%%%%%%%%%%%%%%%%%%%%%%%%%%%%%%%%%%%%%%%%%%%%%%%%%%%%%%%%%%%%%%%%%%%%%%%%%%%%%%%%%%%%%%%%%
\begin{corriges}{
\renewcommand{\carrenoir}{\hspace{-4mm}\rule{1.5mm}{1.5mm}\hspace{2.5mm}}
We give a few programs; the others are in the folder
{\idt book/\allowbreak 07\_for}:
{
\parindent \parindentprog
\parskip 0pt
\nopagebreak\medskip

\carrenoir
\vspace{\neglskip}\mbox{\blc {\mcl from}{\idt \ LibRobot\ }{\mcl import}{\idt \ *}}\par
\nopagebreak
\vspace{0.3\baselineskip}
\vspace{\neglskip}\mbox{\blc {\cmt \#\ Draw\ a\ notched\ frieze}}\par
\nopagebreak
\vspace{0.3\baselineskip}
\vspace{\neglskip}\mbox{\blc {\cmt \#\ Draw\ the\ pattern\ \textquotesingle{}n\_\textquotesingle{}}}\par
\nopagebreak
\vspace{\neglskip}\mbox{\blc {\mcl def}{\idt \ notch():}}\par
\nopagebreak
\vspace{\neglskip}\mbox{\blc \blk\blk{\idt fd();\ qtr();\ fd();\ qtr();\ fd();\ qtl();\ fd()}}\par
\nopagebreak
\vspace{\neglskip}\mbox{\blc \blk\blk{\idt qtl()}}\par
\nopagebreak
\vspace{0.3\baselineskip}
\vspace{\neglskip}\mbox{\blc {\idt gl()}}\par
\nopagebreak
\vspace{\neglskip}\mbox{\blc {\mcl for}{\idt \ i\ }{\mcl in}{\idt \ range(20):}}\par
\nopagebreak
\vspace{\neglskip}\mbox{\blc \blk\blk{\idt notch()}}\par
\nopagebreak
\vspace{\neglskip}\mbox{\blc {\idt st()}}\par

\nopagebreak\medskip }\nopagebreak

{
\parindent \parindentprog
\parskip 0pt
\nopagebreak\medskip

\carrenoir
\vspace{\neglskip}\mbox{\blc {\mcl from}{\idt \ LibRobot\ }{\mcl import}{\idt \ *}}\par
\nopagebreak
\vspace{0.3\baselineskip}
\vspace{\neglskip}\mbox{\blc {\cmt \#\ Draw\ a\ triangular\ grid\ pattern}}\par
\vspace{\neglskip}\mbox{\blc {\cmt \#\ Draw\ 10\ rows\ of\ squares\ side\ by\ side}}\par
\vspace{\neglskip}\mbox{\blc {\cmt \#\ horizontally}}\par
\vspace{\neglskip}\mbox{\blc {\idt qtr()}}\par
\vspace{\neglskip}\mbox{\blc {\mcl for}{\idt \ i\ }{\mcl in}{\idt \ range(1,\ 11):}}\par
\vspace{\neglskip}\mbox{\blc \blk\blk{\cmt \#\ Draw\ a\ row\ of\ i\ squares\ side\ by\ side}}\par
\vspace{\neglskip}\mbox{\blc \blk\blk{\cmt \#\ horizontally,\ starting\ at\ the\ top\ left\ and}}\par
\vspace{\neglskip}\mbox{\blc \blk\blk{\cmt \#\ arriving\ at\ the\ top\ right,\ the\ robot\ facing}}\par
\vspace{\neglskip}\mbox{\blc \blk\blk{\cmt \#\ east\ in\ both\ cases.}}\par
\vspace{\neglskip}\mbox{\blc \blk\blk{\mcl for}{\idt \ j\ }{\mcl in}{\idt \ range(i):}}\par
\vspace{\neglskip}\mbox{\blc \blk\blk\blk\blk{\cmt \#\ Draw\ a\ square,\ starting\ from\ the\ upper}}\par
\vspace{\neglskip}\mbox{\blc \blk\blk\blk\blk{\cmt \#\ left\ corner\ and\ arriving\ at\ the\ upper}}\par
\vspace{\neglskip}\mbox{\blc \blk\blk\blk\blk{\cmt \#\ right\ corner,\ the\ robot\ facing\ east\ in}}\par
\vspace{\neglskip}\mbox{\blc \blk\blk\blk\blk{\cmt \#\ both\ cases.}}\par
\vspace{\neglskip}\mbox{\blc \blk\blk\blk\blk{\mcl for}{\idt \ k\ }{\mcl in}{\idt \ range(4):}}\par
\vspace{\neglskip}\mbox{\blc \blk\blk\blk\blk\blk\blk{\idt fd();\ qtr()}}\par
\vspace{\neglskip}\mbox{\blc \blk\blk\blk\blk{\idt fd()}}\par
\vspace{\neglskip}\mbox{\blc \blk\blk{\cmt \#\ Go\ back,\ without\ drawing,\ to\ the\ start\ of\ the}}\par
\vspace{\neglskip}\mbox{\blc \blk\blk{\cmt \#\ next\ row,\ the\ robot\ facing\ east\ at}}\par
\vspace{\neglskip}\mbox{\blc \blk\blk{\cmt \#\ departure\ and\ on\ arrival.}}\par
\vspace{\neglskip}\mbox{\blc \blk\blk{\idt pu();\ ht()}}\par
\vspace{\neglskip}\mbox{\blc \blk\blk{\mcl for}{\idt \ j\ }{\mcl in}{\idt \ range(i):}}\par
\vspace{\neglskip}\mbox{\blc \blk\blk\blk\blk{\idt fd()}}\par
\vspace{\neglskip}\mbox{\blc \blk\blk{\idt qtl();\ fd();\ qtl();\ pd()}}\par
\nopagebreak
\vspace{\neglskip}\mbox{\blc {\idt st()}}\par

\nopagebreak\medskip }\nopagebreak

{
\parindent \parindentprog
\parskip 0pt
\nopagebreak\medskip

\carrenoir
\vspace{\neglskip}\mbox{\blc {\mcl from}{\idt \ LibRobot\ }{\mcl import}{\idt \ *}}\par
\nopagebreak
\vspace{0.3\baselineskip}
\vspace{\neglskip}\mbox{\blc {\cmt \#\ Draw\ a\ tiling\ of\ squares,\ triangles\ and\ circles}}\par
\vspace{0.3\baselineskip}
\vspace{\neglskip}\mbox{\blc {\cmt \#\ Draw\ a\ square\ with\ sides\ of\ size\ 2,\ the\ robot}}\par
\vspace{\neglskip}\mbox{\blc {\cmt \#\ starting\ from\ the\ lower\ left\ corner\ facing\ east,}}\par
\vspace{\neglskip}\mbox{\blc {\cmt \#\ and\ arriving\ pen\ down\ at\ the\ lower\ right\ corner}}\par
\vspace{\neglskip}\mbox{\blc {\cmt \#\ also\ facing\ east.}}\par
\vspace{\neglskip}\mbox{\blc {\mcl def}{\idt \ square():}}\par
\vspace{\neglskip}\mbox{\blc \blk\blk{\cmt \#\ Draw\ the\ square.}}\par
\vspace{\neglskip}\mbox{\blc \blk\blk{\idt fd();\ fd();\ qtl();\ fd();\ fd();\ qtl()}}\par
\vspace{\neglskip}\mbox{\blc \blk\blk{\idt fd();\ fd();\ qtl();\ fd();\ fd();\ qtl()}}\par
\vspace{\neglskip}\mbox{\blc \blk\blk{\cmt \#\ Advance\ to\ the\ lower\ right\ corner.}}\par
\vspace{\neglskip}\mbox{\blc \blk\blk{\idt fd();\ fd()}}\par
\vspace{0.3\baselineskip}
\vspace{\neglskip}\mbox{\blc {\cmt \#\ Draw\ two\ triangles\ joined\ at\ a\ vertex}}\par
\vspace{\neglskip}\mbox{\blc {\cmt \#\ inscribed\ in\ an\ invisible\ square\ of\ the\ grid\ of}}\par
\vspace{\neglskip}\mbox{\blc {\cmt \#\ side\ 2,\ starting\ from\ the\ lower\ left\ corner\ of\ the\ square}}\par
\vspace{\neglskip}\mbox{\blc {\cmt \#\ facing\ east,\ and\ arriving\ at\ the\ lower\ right}}\par
\vspace{\neglskip}\mbox{\blc {\cmt \#\ corner,\ pen\ down\ and\ facing\ east.}}\par
\vspace{\neglskip}\mbox{\blc {\mcl def}{\idt \ triangles():}}\par
\vspace{\neglskip}\mbox{\blc \blk\blk{\idt pu();\ fd();\ pd()}}\par
\vspace{\neglskip}\mbox{\blc \blk\blk{\idt qtl();\ fd();\ fd();\ tetl();\ fd();\ tetl()}}\par
\vspace{\neglskip}\mbox{\blc \blk\blk{\idt fd();\ fd();\ tetr();\ fd();\ tetl()}}\par
\vspace{\neglskip}\mbox{\blc \blk\blk{\idt pu();\ fd();\ pd()}}\par
\vspace{0.3\baselineskip}
\vspace{\neglskip}\mbox{\blc {\cmt \#\ Draw\ a\ circle\ of\ radius\ 1\ inscribed\ in\ an\ invisible}}\par
\vspace{\neglskip}\mbox{\blc {\cmt \#\ square\ of\ the\ grid\ of\ side\ 2,\ starting\ from\ the}}\par
\vspace{\neglskip}\mbox{\blc {\cmt \#\ lower\ left\ corner\ of\ the\ square\ facing\ east,\ and}}\par
\vspace{\neglskip}\mbox{\blc {\cmt \#\ arriving\ at\ the\ lower\ right\ corner\ of\ the\ square,\ pen}}\par
\vspace{\neglskip}\mbox{\blc {\cmt \#\ down\ and\ facing\ east.}}\par
\vspace{\neglskip}\mbox{\blc {\mcl def}{\idt \ circle():}}\par
\vspace{\neglskip}\mbox{\blc \blk\blk{\idt pu();\ fd();\ pd()}}\par
\vspace{\neglskip}\mbox{\blc \blk\blk{\idt cl();\ cl();\ cl();\ cl()}}\par
\vspace{\neglskip}\mbox{\blc \blk\blk{\idt pu();\ fd();\ pd()}}\par
\vspace{0.3\baselineskip}
\vspace{\neglskip}\mbox{\blc {\cmt \#\ Draw\ a\ row\ made\ of\ a\ pattern\ repeated\ 7}}\par
\vspace{\neglskip}\mbox{\blc {\cmt \#\ times,\ the\ pattern\ being\ made\ of\ a\ square,\ triangles}}\par
\vspace{\neglskip}\mbox{\blc {\cmt \#\ and\ a\ circle\ side\ by\ side,\ the\ robot\ starting\ from\ the}}\par
\vspace{\neglskip}\mbox{\blc {\cmt \#\ lower\ left\ corner\ of\ the\ row\ facing\ east,\ and}}\par
\vspace{\neglskip}\mbox{\blc {\cmt \#\ arriving\ at\ the\ lower\ right\ corner\ also\ facing}}\par
\vspace{\neglskip}\mbox{\blc {\cmt \#\ east.}}\par
\vspace{\neglskip}\mbox{\blc {\mcl def}{\idt \ rowSquare():}}\par
\vspace{\neglskip}\mbox{\blc \blk\blk{\mcl for}{\idt \ i\ }{\mcl in}{\idt \ range(7):}}\par
\vspace{\neglskip}\mbox{\blc \blk\blk\blk\blk{\cmt \#\ The\ robot\ is\ pen\ down,\ facing\ east,\ at\ the\ bottom}}\par
\vspace{\neglskip}\mbox{\blc \blk\blk\blk\blk{\cmt \#\ left\ of\ the\ pattern}}\par
\vspace{\neglskip}\mbox{\blc \blk\blk\blk\blk{\idt square();\ triangles();\ circle()}}\par
\vspace{0.3\baselineskip}
\vspace{\neglskip}\mbox{\blc {\cmt \#\ Like\ rowSquare,\ with\ the\ pattern\ made\ of\ triangles,}}\par
\vspace{\neglskip}\mbox{\blc {\cmt \#\ a\ circle\ and\ a\ square.}}\par
\vspace{\neglskip}\mbox{\blc {\mcl def}{\idt \ rowTriangles():}}\par
\vspace{\neglskip}\mbox{\blc \blk\blk{\mcl for}{\idt \ i\ }{\mcl in}{\idt \ range(7):}}\par
\vspace{\neglskip}\mbox{\blc \blk\blk\blk\blk{\idt triangles();\ circle();\ square()}}\par
\vspace{0.3\baselineskip}
\vspace{\neglskip}\mbox{\blc {\cmt \#\ Like\ rowSquare,\ with\ the\ pattern\ made\ of\ a\ circle,}}\par
\vspace{\neglskip}\mbox{\blc {\cmt \#\ a\ square\ and\ triangles.}}\par
\vspace{\neglskip}\mbox{\blc {\mcl def}{\idt \ rowCircle():}}\par
\vspace{\neglskip}\mbox{\blc \blk\blk{\mcl for}{\idt \ i\ }{\mcl in}{\idt \ range(7):}}\par
\vspace{\neglskip}\mbox{\blc \blk\blk\blk\blk{\idt circle();\ square();\ triangles()}}\par
\vspace{0.3\baselineskip}
\vspace{\neglskip}\mbox{\blc {\cmt \#\ Bring\ the\ robot\ from\ the\ lower\ right\ corner\ of\ a\ row}}\par
\vspace{\neglskip}\mbox{\blc {\cmt \#\ to\ the\ lower\ left\ corner\ of\ the\ next\ row,\ without}}\par
\vspace{\neglskip}\mbox{\blc {\cmt \#\ drawing,\ facing\ east\ and\ pen\ down.}}\par
\vspace{\neglskip}\mbox{\blc {\mcl def}{\idt \ newRow():}}\par
\vspace{\neglskip}\mbox{\blc \blk\blk{\idt pu();\ gl();\ qtr();\ fd();\ fd();\ qtl();\ pd()}}\par
\vspace{0.3\baselineskip}
\vspace{\neglskip}\mbox{\blc {\cmt \#\ Tile\ the\ screen\ with\ 11\ rows\ of\ patterns\ starting\ with}}\par
\vspace{\neglskip}\mbox{\blc {\cmt \#\ a\ square,\ then\ with\ triangles,\ then\ with\ a}}\par
\vspace{\neglskip}\mbox{\blc {\cmt \#\ circle,\ from\ the\ upper\ right\ corner\ to\ the\ lower}}\par
\vspace{\neglskip}\mbox{\blc {\cmt \#\ right\ corner\ of\ the\ screen,\ the\ robot\ facing\ east.}}\par
\vspace{\neglskip}\mbox{\blc {\mcl def}{\idt \ tiling():}}\par
\vspace{\neglskip}\mbox{\blc \blk\blk{\cmt \#\ Draw\ rows\ 1\ to\ 9\ from\ top\ to\ bottom.}}\par
\vspace{\neglskip}\mbox{\blc \blk\blk{\mcl for}{\idt \ j\ }{\mcl in}{\idt \ range(3):}}\par
\vspace{\neglskip}\mbox{\blc \blk\blk\blk\blk{\idt newRow();\ rowSquare()}}\par
\vspace{\neglskip}\mbox{\blc \blk\blk\blk\blk{\idt newRow();\ rowTriangles()}}\par
\vspace{\neglskip}\mbox{\blc \blk\blk\blk\blk{\idt newRow();\ rowCircle()}}\par
\vspace{\neglskip}\mbox{\blc \blk\blk{\cmt \#\ Draw\ rows\ 10\ then\ 11.}}\par
\vspace{\neglskip}\mbox{\blc \blk\blk{\idt newRow();\ rowSquare()}}\par
\vspace{\neglskip}\mbox{\blc \blk\blk{\idt newRow();\ rowTriangles()}}\par
\vspace{0.3\baselineskip}
\vspace{\neglskip}\mbox{\blc {\cmt \#\ Go\ to\ the\ upper\ right\ corner,\ facing\ east.}}\par
\vspace{\neglskip}\mbox{\blc {\idt gt();\ gr();\ qtr()}}\par
\vspace{\neglskip}\mbox{\blc {\cmt \#\ Tile\ the\ screen.}}\par
\nopagebreak
\vspace{\neglskip}\mbox{\blc {\idt tiling();\ st()}}\par

\nopagebreak\medskip }\nopagebreak

\noindent $\Box$\vfill}\end{corriges}

%%%%%%%%%%%%%%%%%%%%%%%%%%%%%%%%%%%%%%%%%%%%%%%%%%%%%%%%%%%%%%%%%%%%%%%%%%%
\noindent \framebox[\textwidth]{
\begin{minipage}{\myparboxwidth}
\setlength{\parindent}{2em}
\small \sf
\vspace{\abovedisplayskip}
\section*{Addenda}

The ``for'' loop is one of the first linguistic notions to have
appeared in computer science, since it was invented at ETH Zürich
 in 1952 by Heinz Rutishauser \cite{rutishauser}.

Comparing the drawings that can be obtained using the ``for''
loop with those that can be made in interactive mode gives a measure
of the power of computers when it comes to performing repetitive tasks
quickly. In this respect the computer is very clearly different
from non-programmable calculators. We also understand
that this computing power has nothing to do with any form
of thought or intelligence, except that of the programmer of
the computer.

In \python, we use a loop {\mcl for} {\idt i} {\mcl in} {\idt range(}$n_1${\idt ,\ }$n_2${\idt ):}
 \ldots\ to repeat
the body of the loop (the indented statements that follow
the colon) $n_2 - n_1$ times, unless $n_2 \leq n_1$, in which case the body of the loop is not
executed. At each iteration, the loop counter {\idt i} takes the
successive values $n_1$, $n_1 + 1$, $n_1 + 2$, \ldots, $n_2 - 1$. In the form {\mcl for}
{\idt i} {\mcl in} {\idt range(}$n_1${\idt ,\ }$n_2${\idt ,\ -1):} \ldots,
{\idt i} takes the successive values $n_1$, $n_1 - 1$, $n_1 - 2$, \ldots, $n_2 + 1$,
unless $n_2 \geq n_1$, in which case the loop is not executed; for example,
{\mcl for} {\idt i} {\mcl in} {\idt range(}$n${\idt ,\ 0,\ -1):} makes {\idt i} take the
values $n$, $n-1$, \ldots, 1 (the third number, here $-1$, is the step
that is added to the counter at each iteration). If the body of the loop modifies the
value of the counter, this has no effect on the following iterations: at each
iteration, the counter takes the next value of the {\idt range} sequence.
When the counter is of no use in the body of the loop, \python\ programmers
often name it `\_', as in
{\mcl for}\ {\idt \_\ }{\mcl in}\ {\idt range(16):}\ \ldots

Using ``for'' loops trains the ability to perceive
when a drawing can be decomposed into a number of identical
sub-drawings (in the case of a simple loop) or similar ones (in the case
of nested loops). This is a daily exercise for the computer scientist,
who must know how to decompose a complex task into simple subtasks,
which sometimes requires an enormous amount of imagination.
It is a first step towards acquiring the inductive
reasoning of sci\-en\-tists who, from specific observations,
create a general hypothesis accounting for these observations,
a hypothesis that remains to be validated in all possible cases, here in
a completely experimental way, by looking at the drawing on
the screen. Mathematicians call such a hypothesis,
induced from experience and not yet proved, a {\em conjecture}.

\vspace{\abovedisplayskip}
\end{minipage}
}

%%%%%%%%%%%%%%%%%%%%%%%%%%%%%%%%%%%%%%%%%%%%%%%%%%%%%%%%%%%%%%%%%%%%%%%%%%%%
%%%%%%%%%%%%%%%%%%%%%%%%%%%%%%%%%%%%%%%%%%%%%%%%%%%%%%%%%%%%%%%%%%%%%%%%%%%
\noindent \framebox[\textwidth]{
\begin{minipage}{\myparboxwidth}
\setlength{\parindent}{2em}
\small \sf
\vspace{\abovedisplayskip}
This inductive reasoning was long confused with reasoning
by mathematical induction, which is essential in mathematics. It took more than a
millennium for mathematicians
(Fermat\footnote{\ \href{https://en.wikipedia.org/wiki/Pierre_de_Fermat}{Pierre de Fermat}, 1601--1665, French mathematician.},
Pascal\footnote{\ \href{https://en.wikipedia.org/wiki/Blaise_Pascal}{Blaise Pascal}, 1623--1662, French mathematician.}
in particular) to fully
understand reasoning by induction (\cite{knuth2} page 17).

For the time being
we are only at the beginnings, the absence of a correct induction
 still being accepted in reasoning
by induction: if the drawing is right
for the basic pattern and for
the first repetition of the pattern, then it is deemed correct, the only
remaining problem, fitting the drawing into the window, being solved by simple
experimentation.

The notion of {\em invariant} (that is, for our \robot, the
description of its possible positions and orientations at the beginning of the
body of the loops) is due to Turing%
\footnote{\ \href{https://en.wikipedia.org/wiki/Alan_Turing}{Alan Turing}, English computer scientist.} in 1949; it was forgotten, then rediscovered by
Naur%
\footnote{\ \href{https://en.wikipedia.org/wiki/Peter_Naur}{Peter Naur}, Danish computer scientist.} in 1966 and Floyd%
\footnote{\ \href{https://en.wikipedia.org/wiki/Robert_W._Floyd}{Robert Floyd}, American computer scientist.} in 1967.
It relies
on reasoning by induction: if the invariant is true when we
enter the loop for the first time (that is, the
position and orientation of the \robot\ conform to the description
given in the invariant) and if, assuming that the invariant is true at the
beginning of the execution of the body of the loop, it is still true after
the execution of this loop body, then the invariant is always true at the
beginning of each iteration of the loop and is still true when
we exit the loop. The notion applies equally well to procedures: we
can use an
{\em entry invariant} describing the positions and orientations of the
\robot\ before the call and an {\em exit invariant} describing
these positions and orientations after the call, so that the body of the
procedure can be understood without having to look at all the places
where it is called to find out under which conditions the procedure is called.

\vspace{\abovedisplayskip}
\end{minipage}
}

%%%%%%%%%%%%%%%%%%%%%%%%%%%%%%%%%%%%%%%%%%%%%%%%%%%%%%%%%%%%%%%%%%%%%%%%%%%%
%%%%%%%%%%%%%%%%%%%%%%%%%%%%%%%%%%%%%%%%%%%%%%%%%%%%%%%%%%%%%%%%%%%%%%%%%%%
\noindent \framebox[\textwidth]{
\begin{minipage}{\myparboxwidth}
\setlength{\parindent}{2em}
\small \sf
\vspace{\abovedisplayskip}
The computer can be used to manipulate programs, in
particular to discover automatically invariants of simple form
at each point of a program. This is the purpose of the
{\em abstract interpretation} of programs
\cite{CousotIA}. For the \robot, for example, a simple form
of invariant is the set of its possible orientations,
ignoring the position of its pen, its coordinates on
the screen, etc. The possible orientations of the concrete \robot\ can be
computed by executing the program with an abstract \robot\ that
executes the commands of the concrete \robot\ without moving
but simply changing its orientation. For this abstract \robot,
the command {\idt fd} does nothing since it does not change the orientation of the
concrete \robot; the command {\idt cr} has the same effect as
{\idt qtr}; etc. During the abstract execution, the possible orientations
of the abstract \robot\ are written as comments
at each point of the program. The abstract execution of the program
stops once the results have stabilized. For example, for the program {\idt FriezeOfN\_2} (page
\pageref{prog:FriseDeN_2}), we find:

{
\parindent \parindentprog
\parskip 0pt
\nopagebreak\medskip

\carrenoir
\vspace{\neglskip}\mbox{\blc {\mcl from}{\idt \ LibRobot\ }{\mcl import}{\idt \ *}}\par
\nopagebreak
\vspace{0.3\baselineskip}
\vspace{\neglskip}\mbox{\blc {\cmt \#\ Draw\ a\ frieze\ of\ N\ with\ a\ \char34{}for\char34{}\ loop}}\par
\vspace{\neglskip}\mbox{\blc {\cmt \#\ North}}\par
\vspace{\neglskip}\mbox{\blc {\idt gl()}}\par
\vspace{\neglskip}\mbox{\blc {\cmt \#\ North}}\par
\vspace{\neglskip}\mbox{\blc {\mcl for}{\idt \ j\ }{\mcl in}{\idt \ range(16):}}\par
\vspace{\neglskip}\mbox{\blc \blk\blk{\idt fd();\ fd()\ \ }{\cmt \#\ North\ \ \#\ North}}\par
\vspace{\neglskip}\mbox{\blc \blk\blk{\idt ht();\ cl()\ \ }{\cmt \#\ North\ \ \#\ South}}\par
\vspace{\neglskip}\mbox{\blc \blk\blk{\idt cr();\ ht()\ \ }{\cmt \#\ East\ \ \#\ South}}\par
\vspace{\neglskip}\mbox{\blc {\cmt \#\ North}}\par
\vspace{\neglskip}\mbox{\blc {\cmt \#\ North}}\par
\vspace{\neglskip}\mbox{\blc {\idt st()}}\par
\nopagebreak
\vspace{\neglskip}\mbox{\blc {\cmt \#\ North}}\par

\nopagebreak\medskip }\nopagebreak


The rule of signs is an example of abstract interpretation used
in mathematics. The sign of an expression can be
computed using the sign of the variables, which replaces their values,
thanks to abstract operations such as $\oplus + \oplus = \oplus$,
$\oplus - \ominus = \oplus$, $\oplus \times \ominus = \ominus$, etc. The price
to pay for this simplification is of course the imprecision of the
results, as for $\oplus + \ominus$.
\vspace{\abovedisplayskip}
\end{minipage}
}
%%%%%%%%%%%%%%%%%%%%%%%%%%%%%%%%%%%%%%%%%%%%%%%%%%%%%%%%%%%%%%%%%%%%%%%%%%%%
